\section{Preentrenamiento: la experiencia reproducible de OLMo 2}

\subsection{Objetivo: producir un modelo base muy competitivo dentro de un costo controlable}
El expositor usa OLMo 2 para discutir una realidad clave: los equipos de investigación no siempre cuentan con poder de cómputo de nivel frontier, pero aun así pueden, mediante una mejor gobernanza de datos y detalles de entrenamiento, producir una línea base sólida en las escalas de 7B/13B.

\begin{importantbox}{El valor de OLMo 2 está en la ``mejora interpretable''}
En comparación con solo publicar las curvas de resultados, la serie OLMo enfatiza más la publicación abierta de la cadena de entrenamiento, lo cual permite a la comunidad responder de manera sistemática:
\begin{itemize}
    \item qué detalles de entrenamiento realmente contribuyeron al desempeño;
    \item qué modificaciones solo son efectivas en puntos de referencia específicos;
    \item bajo distintos presupuestos, qué conjunto de configuraciones es más robusto.
\end{itemize}
\end{importantbox}

\begin{figure}[H]
\centering
\includegraphics[width=0.9\textwidth]{frame-03-olmo-scaling.jpg}
\caption{Video aproximadamente en 00:18:30: curvas de entrenamiento de OLMo y comparación de escalas, con énfasis en la optimización de la eficiencia bajo presupuesto restringido}
\end{figure}

\subsection{Los factores que realmente influyen en el modelo base de razonamiento durante el preentrenamiento}
\begin{knowledgebox}{Cuatro tipos de variables que se pasan por alto con más facilidad que la ``escala de parámetros''}
\begin{enumerate}
    \item la deduplicación de datos, el control de contaminación y la estratificación por calidad;
    \item la proporción de muestreo por dominio (código, matemáticas, texto de conocimiento, diálogo);
    \item la tasa de aprendizaje y la programación de warmup/cooldown;
    \item el tokenizer y la estrategia de ventana de contexto.
\end{enumerate}
\end{knowledgebox}

El expositor enfatiza que la capacidad de razonamiento ya queda parcialmente determinada en la etapa de preentrenamiento, y que el posentrenamiento solo puede reordenar y reforzar dentro de los límites de la capacidad ya existente. Esto explica por qué muchas técnicas de posentrenamiento producen resultados muy distintos según el modelo base.

\begin{warningbox}{Las métricas del preentrenamiento no corresponden una a una con la capacidad de razonamiento en tareas posteriores}
Un loss de modelado de lenguaje más bajo no necesariamente produce un razonamiento en cadena más sólido. Si el corpus de entrenamiento tiene señales insuficientes de estructura de razonamiento, el modelo puede ser hábil para continuar textos con fluidez, pero no para la inferencia con restricciones de varios pasos.
\end{warningbox}

\subsection{Resumen de la sección}
La lección de OLMo 2 es que, en un entorno de investigación, una receta de preentrenamiento abierta e interpretable constituye en sí misma un resultado, y determina el límite superior y la reproducibilidad de todas las mejoras de reasoning posteriores.


